% !TeX root = ../../Book2.tex
% 纲要标题：### 4.4 线性算子的分类
% 纲要位置：计划纲要.md，第 936 行。

\section{线性算子的分类}
\label{b2:sec:0404}

% 纲要要点（供目录核对）：
% * k[t] 模与算子相似
% * 循环分解及有理标准形
% * 最小多项式与特征多项式

设 \(k\) 为任意域，\(V\) 为有限维 \(k\) 向量空间，\(T:V\to V\) 为线性算子。第一章把 \(T\) 写成多项式作用，第二章将其定义为 \(k[t]\) 模 \(V_T\)。现在 \(k[t]\) 是主理想整环，模的分类定理已经可用；所得标准形适用于有限域，也不要求特征多项式在 \(k\) 中分裂。

\subsection{多项式环模与算子相似}
\label{b2:subsec:0404:01}

\begin{proposition}\label{b2:prop:operator-similarity-modules}
设 \(V,W\) 是域 \(k\) 上的有限维向量空间，\(T:V\to V\)、\(S:W\to W\) 是线性算子；\(V_T,W_S\) 表示分别由 \(t\cdot v=Tv\)、\(t\cdot w=Sw\) 得到的 \(k[t]\) 模。映射 \(u:V_T\to W_S\) 是 \(k[t]\) 模同态，当且仅当 \(u\) 是 \(k\) 线性映射并满足 \(uT=Su\)。因此两个同维算子的矩阵相似，当且仅当对应 \(k[t]\) 模同构。
\end{proposition}
\begin{proof}
同态的判据已由\cref{b2:prop:operator-module-intertwiners}证明：保持常数标量给出 \(k\) 线性，保持 \(t\) 的作用给出 \(uT=Su\)，而这个等式又保证保持所有多项式的作用。若 \(u\) 双射，分别取 \(V,W\) 的基，以 \(A,B\) 表示 \(T,S\)，以 \(P\) 表示同构 \(u\) 的矩阵，则 \(P\) 可逆，等式化为 \(PA=BP\)，即 \(B=PAP^{-1}\)。令 \(Q=P^{-1}\)，便有 \(B=Q^{-1}AQ\)，与第一章的相似约定相同；反过来，相似矩阵所给的可逆映射也定义模同构。
\end{proof}

取 \(k^n\) 的标准基 \(e_1,\ldots,e_n\)，在算子模中有
\[
te_j=Ae_j=\sum_i a_{ij}e_i.
\]
若先以这些符号生成自由模 \(k[t]^n\)，再加入关系 \(te_j-\sum_i a_{ij}e_i=0\)，这 \(n\) 条关系按列排成的矩阵正是 \(tI_n-A\)。要得到 \(V_A\) 的展示，还须证明它们生成全部关系。下面的求值映射将完成这一步。这里对 \(A\) 的分类是 \(k\) 上的相似，对 \(tI-A\) 的化简则是 \(k[t]\) 上左右独立的幺模等价。

\begin{proposition}\label{b2:prop:operator-polynomial-matrix}
设 \(k\) 是域，\(A\in M_n(k)\)，并把 \(k^n\) 配上 \(t\cdot v=Av\) 的 \(k[t]\) 模结构，记为 \(V_A\)，则有 \(k[t]\) 模同构
\[
V_A\cong\operatorname{coker}(tI_n-A).
\]
具体地，映射
\[
\varepsilon:k[t]^n\longrightarrow V_A,\qquad
\sum_{j=0}^d t^jv_j\longmapsto\sum_{j=0}^d A^jv_j
\]
是满模同态，核为 \((tI_n-A)k[t]^n\)。
\end{proposition}
\begin{proof}
\(\varepsilon\) 保持常数标量，并满足 \(\varepsilon(tF)=A\varepsilon(F)\)，故对多项式有限求和可知它 \(k[t]\) 线性。任意向量是相应常数向量多项式的像，所以它满射。由于
\[
\varepsilon\bigl((tI_n-A)G\bigr)
=A\varepsilon(G)-A\varepsilon(G)=0,
\]
所述关系子模包含在核中。

反过来，模掉这些关系后，\(tv\equiv Av\)，反复替换便有 \(t^jv\equiv A^jv\)。因此每个多项式向量 \(F=\sum_jt^jv_j\) 都应与常数向量 \(\varepsilon(F)\) 同余。这个替换过程由矩阵恒等式
\[
t^jI_n-A^j=(tI_n-A)\sum_{h=0}^{j-1}t^{j-1-h}A^h
\qquad(j\ge1)
\]
严格给出；右边展开后中间项逐项抵消，只剩左边。于是
\[
F-\varepsilon(F)
=\sum_{j\ge1}(t^jI_n-A^j)v_j
\in(tI_n-A)k[t]^n.
\]
这里把 \(\varepsilon(F)\) 看成常数向量多项式。若 \(\varepsilon(F)=0\)，便得 \(F\) 属于所述关系子模，证明核的反向包含。最后应用\cref{b2:thm:module-first-isomorphism}得到同构。
\end{proof}

\subsection{循环分解与有理标准形}
\label{b2:subsec:0404:02}

一个 \(k\) 基当然也是 \(V_T\) 的 \(k[t]\) 模生成集，所以它有限生成。对每个 \(v\)，向量 \(v,Tv,\ldots,T^nv\) 线性相关，某个非零多项式因而零化 \(v\)。所以 \(V_T\) 是挠 \(k[t]\) 模，\cref{b2:thm:pid-module-structure}中的自由部分为零。

\begin{theorem}[有理标准形]\label{b2:thm:rational-canonical-form}
设 \(V\) 是域 \(k\) 上的有限维向量空间，\(T\in\operatorname{End}_k(V)\)，并以 \(t\cdot v=Tv\) 将 \(V\) 看成 \(k[t]\) 模 \(V_T\)。存在唯一一列首一非常数多项式
\[
f_1\mid f_2\mid\cdots\mid f_s
\]
使 \(V_T\cong\bigoplus_{i=1}^s k[t]/(f_i)\)。若 \(V=0\)，列表为空。对每个
\[
f_i=t^{d_i}+a_{d_i-1}t^{d_i-1}+\cdots+a_0,
\]
记其友矩阵为
\[
C(f_i)=\begin{pmatrix}
0&0&\cdots&0&-a_0\\
1&0&\cdots&0&-a_1\\
0&1&\cdots&0&-a_2\\
\vdots&\vdots&\ddots&\vdots&\vdots\\
0&0&\cdots&1&-a_{d_i-1}
\end{pmatrix}.
\]
当 \(d_i=1\) 时此式解释为单项矩阵 \((-a_0)\)。\(V\) 有一组基，使 \(T\) 的矩阵为 \(C(f_1)\oplus\cdots\oplus C(f_s)\)，称为有理标准形。两个算子相似，当且仅当这列多项式相同。
\end{theorem}
\begin{proof}
\(k[t]\) 为主理想整环，且 \(V_T\) 有限生成并为挠模，由\cref{b2:thm:pid-module-structure}得到循环商分解。\cref{b2:thm:pid-elementary-divisors}保证非单位因子在相伴意义下唯一；取首一代表，得到唯一的整除链。

带余除法说明 \(k[t]/(f_i)\) 中 \(1,t,\ldots,t^{d_i-1}\) 的陪集是 \(k\) 基。乘 \(t\) 把前 \(d_i-1\) 项送到下一项，而最后一项满足
\[
t^{d_i}=-a_0-a_1t-\cdots-a_{d_i-1}t^{d_i-1}\pmod{f_i},
\]
所以该算子的矩阵正是 \(C(f_i)\)。将每个循环商的基通过模同构送入 \(V\)，再把它们合并，得到所需基与分块矩阵。最后，\cref{b2:prop:operator-similarity-modules}把模同构与矩阵相似等同起来；分类唯一性便给出末项判据。
\end{proof}

每个循环商生成元在 \(V\) 中的像为 \(v_i\)，相应子空间具有基 \(v_i,Tv_i,\ldots,T^{d_i-1}v_i\)。计算 \(tI-A\) 的 Smith 形，可以同时求这些关系多项式；若追踪行变换，\cref{b2:prop:smith-cokernel}还给出可用于恢复循环向量的原坐标。

具体地，设 \(A\in M_n(k)\)，并已求得
\[
P(t)(tI_n-A)Q(t)=\operatorname{diag}(g_1,\ldots,g_n),
\qquad P(t),Q(t)\in\operatorname{GL}_n(k[t]),
\]
其中 \(g_j\) 首一且成整除链。对每个非常数因子 \(g_j\)，Smith 余核的第 \(j\) 个标准向量 \(e_j\) 生成循环商 \(k[t]/(g_j)\)。由\cref{b2:prop:smith-cokernel}，\(P(t)^{-1}\) 把它送回原展示的目标坐标，得到代表 \(P(t)^{-1}e_j\)。这个代表仍是形式多项式向量；要得到原空间 \(k^n\) 中的向量，还须由\cref{b2:prop:operator-polynomial-matrix}中的 \(\varepsilon\) 把多项式作用解释成 \(A\) 的迭代。因此
\[
v_j=\varepsilon\bigl(P(t)^{-1}e_j\bigr)\in k^n.
\]
整个对应可写为
\[
\begin{aligned}
\operatorname{coker}\operatorname{diag}(g_1,\ldots,g_n)
&\xrightarrow{\ P(t)^{-1}\ }\operatorname{coker}(tI_n-A)
\xrightarrow{\ \overline\varepsilon\ }V_A,\\
[e_j]&\longmapsto[P(t)^{-1}e_j]\longmapsto v_j.
\end{aligned}
\]
这里 \(\overline\varepsilon\) 是 \(\varepsilon\) 在余核上诱导的同构。若 \(P(t)^{-1}e_j=\sum_\ell t^\ell w_\ell\)，则 \(v_j=\sum_\ell A^\ell w_\ell\)。令 \(d_j=\deg g_j\)，按因子次序将各组列向量
\(v_j,Av_j,\ldots,A^{d_j-1}v_j\) 合成矩阵 \(S\)。上述商模同构说明这些向量组成 \(k^n\) 的基，故 \(S\) 可逆，且
\[
S^{-1}AS=\bigoplus_{\deg g_j>0}C(g_j).
\]
单位因子对应零商模，不贡献基向量。这就把关系矩阵的幺模变换明确转成了原算子的相似换基。

\ProseParagraphBreak
例如在 \(\mathbb F_2\) 上，\(p=t^3+t+1\) 在 \(0,1\) 处都非零，三次多项式若可约必有一次因子，故 \(p\) 不可约。矩阵
\[
C(p)=\begin{pmatrix}0&0&1\\1&0&1\\0&1&0\end{pmatrix}
\]
本身就是一个有理标准形，标准向量 \(e_1\) 的前三个迭代 \(e_1,C(p)e_1,C(p)^2e_1\) 构成基。它在 \(\mathbb F_2\) 中没有特征值，却完全处在本定理的范围内。“有理”在这里不要求系数域等于 \(\mathbb Q\)，而是指无须为取得特征值而扩大原域。

\subsection{最小多项式与特征多项式}
\label{b2:subsec:0404:03}

记 \(\chi_T(t)=\det(tI-A)\) 为特征多项式，\(\mu_T\) 为首一最小多项式。相似矩阵给出相同的特征多项式；最小多项式记录全部向量共同满足的最低次数首一关系。它们在循环分解中分别体现为乘积与最大的关系因子。下面直接计算友矩阵的行列式，因为本节随后要由循环分解证明 Cayley--Hamilton 定理，不能预先用它认定这个行列式。

\begin{lemma}\label{b2:lem:companion-characteristic}
设 \(k\) 是域，\(f=t^d+a_{d-1}t^{d-1}+\cdots+a_0\in k[t]\) 是首一多项式，\(d\ge1\)，则其友矩阵 \(C(f)\in M_d(k)\) 满足
\[
\det(tI_d-C(f))=f(t).
\]
\end{lemma}
\begin{proof}
\(d=1\) 时等式为 \(t+a_0=f\)。一般情形中，\(tI_d-C(f)\) 的前 \(d-1\) 列只有对角项 \(t\) 与下次对角项 \(-1\)，最后一列为 \((a_0,a_1,\ldots,a_{d-2},t+a_{d-1})^{\mathsf T}\)。按最后一列展开。删除第 \(i\) 行后，其余前 \(d-1\) 列的行列式为 \(t^{i-1}(-1)^{d-i}\)：左上三角块贡献 \(t^{i-1}\)，右下三角块贡献 \((-1)^{d-i}\)。乘代数余子式符号 \((-1)^{i+d}\) 后，系数成为 \(t^{i-1}\)。所以展开式为
\[
a_0+a_1t+\cdots+a_{d-2}t^{d-2}+(t+a_{d-1})t^{d-1}=f(t).
\]
\end{proof}

\begin{proposition}\label{b2:prop:operator-characteristic-minimal}
设 \(V\ne0\) 是域 \(k\) 上的有限维向量空间，\(T\in\operatorname{End}_k(V)\)，且 \(V_T\) 的不变因子为 \(f_1\mid\cdots\mid f_s\)，则
\[
\chi_T=\prod_{i=1}^s f_i,\qquad \mu_T=f_s,
\qquad \dim_kV=\sum_{i=1}^s\deg f_i.
\]
特别地，\(\mu_T\mid\chi_T\) 且 \(\chi_T(T)=0\)。零空间上约定 \(\chi_T=\mu_T=1\)，维数为零。
\end{proposition}
\begin{proof}
在有理标准形的基中，特征矩阵分块对角。行列式为各块行列式的乘积，由\cref{b2:lem:companion-characteristic}得到 \(\chi_T=\prod_i f_i\)。各块具有 \(\deg f_i\) 个基向量，也给出维数公式。

在 \(k[t]/(f_i)\) 中，一个多项式 \(g\) 零化全模，当且仅当它零化 \(1+(f_i)\)，即 \(f_i\mid g\)。因此 \(g(T)=0\) 等价于每个 \(f_i\mid g\)，由整除链又等价于 \(f_s\mid g\)。首一最小关系因而为 \(f_s\)，并且整除乘积 \(\chi_T\)，从而 \(\chi_T(T)=0\)。
\end{proof}

这里也给出了\term[Cayley-Hamilton-dingli]{Cayley--Hamilton 定理}[Cayley--Hamilton theorem]的证明：每个方阵满足自己的特征多项式。\footnote{凯莱（Arthur Cayley，1821-1895），英国数学家，发展矩阵理论与不变量理论。哈密顿（William Rowan Hamilton，1805-1865），爱尔兰数学家及物理学家，创立四元数。}第一章的\cref{b2:prop:minimal-polynomial-divisibility}曾利用 \(I,T,\ldots,T^{n^2}\) 的线性相关性给出最小多项式存在的粗界。现在 \(\mu_T\mid\chi_T\) 且 \(\deg\chi_T=n\)，便将它改进为 \(\deg \mu_T\le n\)。证明所需的挠性来自有限维性，分类过程无须预设 Cayley--Hamilton 定理。

\ProseParagraphBreak
结合\cref{b2:prop:operator-polynomial-matrix}，\(tI-A\) 的非单位 Smith 因子恰为 \(f_1,\ldots,f_s\)。由于 \(\det(tI-A)\) 为次数 \(n\) 的首一多项式，它不为零，所以 Smith 形没有零对角元；其余 \(n-s\) 项是单位，归一化后为 \(1\)。于是
\[
P(t)(tI-A)Q(t)=\operatorname{diag}(1,\ldots,1,f_1,\ldots,f_s)
\]
把特征多项式、最小多项式与原矩阵的子式理想连接起来。这个矩阵恒等式也是下一节讨论基域扩张时保留分类数据的依据。

\ProseParagraphBreak
特征多项式给出全部不变因子的乘积，最小多项式只给出最后一项，即使再给定秩，也通常不足以决定相似类型。\cref{b2:prop:insufficient-similarity-invariants}中的八维算子现在可以写成两列不变因子
\[
(t,t^3,t^4),\qquad (t^2,t^2,t^4).
\]
两列的乘积都是 \(t^8\)，最后一项都是 \(t^4\)。每个 \(k[t]/(t^d)\) 中乘 \(t\) 的核均为一维，所以两个算子的核维数都为三、秩都为五。但两列整除链不同，故算子不相似；平方算子的核维数分别为五与六，能检测到乘积、末项和秩共同遗漏的信息。
